\chapter{Mục tiêu và phạm vi}

\newpage

Đồ án số 2 yêu cầu đi vào tầng phân trang của xv6: đọc và sửa bảng trang để tăng tốc một system call và để phát hiện những trang đã được truy cập. Cụ thể có ba bài tập.

\begin{enumerate}[itemsep=2.5pt, topsep=2.5pt, leftmargin=1.5cm]
    \item \textbf{Mở rộng \texttt{sysinfo}} của đồ án 1 để trả thêm số byte bộ nhớ vật lý còn trống.
    \item \textbf{In bảng trang}: viết \texttt{vmprint()} trong kernel, in cây bảng trang theo đúng định dạng của đề, và giải thích ý nghĩa cùng quyền truy cập của từng trang lá.
    \item \textbf{Phát hiện trang được truy cập}: cài đặt system call \texttt{pgaccess()} đọc bit \texttt{PTE\_A} của một dải trang và trả về bitmask cho user space.
\end{enumerate}

Cả ba bài đều nằm trong cùng một cơ chế: bảng trang Sv39 của RISC-V. Vì vậy chương 2 trình bày trước cách xv6 dịch địa chỉ, cách \texttt{walk()} đi xuống ba cấp và ý nghĩa các bit trong PTE; ba chương sau mỗi chương giải quyết một bài tập, kèm kết quả chạy thật trên QEMU.

Nhóm làm việc trên nhánh \texttt{pgtbl} của bộ \texttt{xv6-labs-2024}. Đề bài không yêu cầu một mục khảo sát mã nguồn riêng như đồ án 1, nên báo cáo chỉ giữ phần lý thuyết tối thiểu cần cho ba bài tập.

\vspace{0.5cm}

\begin{table}[H]
\centering
\renewcommand{\arraystretch}{1.3}
\setlength{\tabcolsep}{8pt}
\begin{tabular}{|c|p{9cm}|c|}
\hline
\rowcolor[rgb]{0.8, 0.9, 1.0}
\textbf{Yêu cầu} & \textbf{Nội dung} & \textbf{Mục tương ứng} \\ \hline
2.1 & Cơ chế phân trang Sv39, \texttt{walk()}, \texttt{mappages()}, bố cục \texttt{memlayout.h} & 2.1--2.4 \\ \hline
3 & Số byte bộ nhớ trống trong \texttt{sysinfo} & 3.1--3.4 \\ \hline
4 & \texttt{vmprint()}, mốc gọi trong \texttt{exec()} & 4.1--4.3 \\ \hline
4 & Mở rộng: in theo yêu cầu bằng system call \texttt{vmprint()} & 4.4 \\ \hline
4 & Giải thích từng trang lá và quyền truy cập & 4.6 \\ \hline
5 & \texttt{pgaccess()}, bit \texttt{PTE\_A}, chương trình kiểm thử & 5.1--5.5 \\ \hline
\end{tabular}
\caption{Đối chiếu yêu cầu của đề và các mục trong báo cáo}
\end{table}
